\begin{lemma}[The positive part of a constant minus an $L^2(\gamma_n)$ function is $L^2(\gamma_n)$]
Fix $n \in \mathbb{N}_0$ and let $\gamma_n$ be the Gaussian measure of
\noderef{GaussMeasure}.  Let $F \in L^2(\gamma_n)$ and let $a \in \mathbb{R}$ be
any real number.  Then the positive part of $a - F$,
\[
  (a - F)^{+}(x) = \max\bigl(a - F(x),\, 0\bigr),
\]
again lies in $L^2(\gamma_n)$.

All three memberships here are with respect to the Gaussian measure $\gamma_n$:
$F$ is square $\gamma_n$-integrable and so is the conclusion's function, the
integrals being taken against $\gamma_n$ and not against Lebesgue measure.  Note
also that no sign condition on $a$ is assumed or needed --- $a$ ranges over all
of $\mathbb{R}$, and the argument below uses nothing about whether $a$ is
positive.
\end{lemma}
\begin{proof}
Step 1 (the constant is square integrable).  A constant function is square
integrable for any finite measure, and $\gamma_n$ is a probability measure by
\noderef{GaussIsProb}, hence finite.  So $x \mapsto a$ lies in $L^2(\gamma_n)$.

Step 2 ($a - F$ is square integrable).  $L^2(\gamma_n)$ is a vector space, so
the difference of the constant $a$ of Step 1 and the given $F \in L^2(\gamma_n)$
lies in $L^2(\gamma_n)$:
\[
  x \longmapsto a - F(x) \quad \in L^2(\gamma_n). \tag{D}
\]

Step 3 (pointwise domination of the positive part).  For every
$x \in \mathbb{R}^n$,
\[
  \bigl| \max(a - F(x), 0) \bigr| \leq \bigl| a - F(x) \bigr| . \tag{P}
\]
To see this, write $t = a - F(x)$ and note first that $\max(t,0) \geq 0$, so the
left-hand side is $\max(t,0)$ itself.  Now distinguish two cases, which together
are exhaustive.  If $t \leq 0$ then $\max(t,0) = 0$, and the right-hand side
$|t|$ is nonnegative, so the inequality reads $0 \leq |t|$ and holds.  If
$t \geq 0$ then $\max(t,0) = t$, and $t \leq |t|$ holds because a real number is
at most its own absolute value.  In both cases (P) follows.  This is the only
step that touches the shape of the positive part, and it is where the absence of
any hypothesis on $a$ shows: the comparison is between $\max(t,0)$ and $|t|$ for
an arbitrary real $t$, with no reference to how $t$ arose.

Step 4 (measurability of the positive part).  The function $F$ is
$\gamma_n$-almost-everywhere strongly measurable, being an element of
$L^2(\gamma_n)$, hence so is $x \mapsto a - F(x)$ by Step 2.  Taking the
pointwise maximum with the constant $0$ preserves almost-everywhere
measurability, since the maximum of two almost-everywhere measurable real
functions is almost-everywhere measurable, and for real-valued functions
almost-everywhere measurability and almost-everywhere strong measurability
agree.  So $x \mapsto \max(a - F(x), 0)$ is $\gamma_n$-almost-everywhere
strongly measurable.

Step 5 (conclusion).  A function that is $\gamma_n$-almost-everywhere strongly
measurable and dominated in absolute value, $\gamma_n$-almost everywhere, by a
function of $L^2(\gamma_n)$ is itself in $L^2(\gamma_n)$.  Step 4 supplies the
measurability, (P) supplies the domination --- indeed at every point, so in
particular $\gamma_n$-almost everywhere --- and (D) says the dominating function
$x \mapsto a - F(x)$ is in $L^2(\gamma_n)$.  Hence
$x \mapsto \max(a - F(x), 0)$ lies in $L^2(\gamma_n)$, as asserted.
\end{proof}
